\chapter{Solution proposée}
\label{ch:solution}

\section{Vision produit}
\label{sec:vision-produit}

La passerelle est un middleware configurable. Elle ne décide pas a priori des dispositifs supportés. Elle expose une mécanique de plug-in qui permet de déclarer, en JSON, comment interpréter une nouvelle source. La philosophie produit se résume en trois questions et deux phases.

\begin{table}[h]
  \centering
  \caption{Philosophie produit : trois questions, deux phases.}
  \label{tab:philosophie}
  \begin{tabularx}{\linewidth}{lXl}
    \toprule
    \bridgeheader Phase & Question & Responsable \\
    \midrule
    Handshake technique & Q1. Quel canal de communication ? (TCP, RS-232, BLE, fichier, MQTT) & Développeur \\
    Handshake technique & Q2. Quelle langue ? (HL7v2, JSON, XML, ASCII MEDIBUS) & Développeur \\
    Configuration métier & Q3. Quelle cible et quel mapping ? (HAPI FHIR, codes LOINC, identité patient) & Admin via dashboard \\
    \bottomrule
  \end{tabularx}
\end{table}

Les deux premières questions sont résolues par le développeur via un profil JSON livré avec la passerelle. La troisième est exposée à l'administrateur via le dashboard. Cette séparation est un principe directeur du produit.

\section{Cas d'usage prioritaires}
\label{sec:use-cases}

Six cas d'usage cadrent la démo.

\begin{table}[h]
  \centering
  \caption{Use cases couverts par la démo.}
  \label{tab:use-cases}
  \begin{tabularx}{\linewidth}{llX}
    \toprule
    \bridgeheader ID & Acteur & Description \\
    \midrule
    UC1 & Moniteur Philips MX450 & Capter une trame HL7 v2 ORU\textasciicircum{}R01 over MLLP, extraire SpO\textsubscript{2} et FC, publier en FHIR sur HAPI. \\
    UC2 & Glucomètre Masimo Radical-7 & Recevoir un message MQTT JSON, mapper la glycémie en LOINC 2339-0, publier en FHIR. \\
    UC3 & ECG Welch Allyn CP150 & Détecter un fichier SCP-ECG déposé dans un volume, parser, publier une étude ECG (LOINC 11524-6). \\
    UC4 & Pompe Dräger Evita V500 & Lire une trame ASCII MEDIBUS sur RS-232 simulé, mapper FiO\textsubscript{2}, PEEP, Vt, publier en FHIR. \\
    UC5 & Auteur du plugin & Ajouter un nouveau dispositif en déposant un plugin sans redémarrage du cœur. \\
    UC6 & Technicien biomédical & Recevoir une alerte mail technique sur valeur hors plage de plausibilité, sans donnée patient. \\
    \bottomrule
  \end{tabularx}
\end{table}

\section{Architecture cible : quatre vues}
\label{sec:architecture-cible}

L'architecture est documentée selon les quatre vues classiques de la pyramide d'architecture.

\subsection{Vue métier}

La passerelle sert deux objectifs métier. Côté clinique, elle alimente le DPI en données vitales fiables et horodatées, support à la décision soignante. Côté OT, elle donne aux techniciens biomédicaux une vision temps réel du parc connecté, sans accès aux données cliniques.

\subsection{Vue fonctionnelle}

La chaîne de traitement est linéaire : capture, parsing, mapping sémantique, transport sécurisé, publication. Chaque étape est un point d'extension.

\begin{figure}[h]
  \centering
  \resizebox{\linewidth}{!}{% =====================================================================
% figures/fig-pipeline-4couches.tex
% Pipeline fonctionnel : Adapter -> Parser -> Mapper -> Transport -> HAPI
% Inputs varies a gauche, sortie FHIR a droite. Bandeau plugins au-dessus.
% =====================================================================
\begin{tikzpicture}[
  font=\sffamily,
  every node/.style={font=\small\sffamily},
  >={Stealth[length=2.5mm]},
  thick,
  layer/.style={
    rectangle, rounded corners=3pt,
    draw=bridgeblue, fill=bridgeblue!10, text=bridgeblue,
    minimum width=2.4cm, minimum height=4.2cm,
    align=center,
    font=\small\bfseries\sffamily
  },
  altlayer/.style={
    rectangle, rounded corners=3pt,
    draw=bridgeteal, fill=bridgeteal!10, text=bridgeteal,
    minimum width=2.4cm, minimum height=4.2cm,
    align=center,
    font=\small\bfseries\sffamily
  },
  outbox/.style={
    rectangle, rounded corners=3pt,
    draw=bridgegreen, fill=bridgegreen!15, text=bridgegreen,
    minimum width=2.0cm, minimum height=1.0cm,
    align=center,
    font=\small\bfseries\sffamily
  },
  inputlbl/.style={
    anchor=east, color=bridgegray,
    font=\footnotesize\itshape\sffamily,
    inner sep=2pt
  },
  caption/.style={
    align=center, font=\scriptsize\itshape\sffamily,
    color=bridgegray, inner sep=1pt
  },
  flowarrow/.style={-{Stealth[length=3mm]}, thick, draw=bridgeblue}
]

% =====================================================================
% Etiquettes des entrees (a gauche)
% =====================================================================
\node[inputlbl] (lblTcp) at (-0.2, 4.5) {TCP / HL7v2};
\node[inputlbl] (lblBle) at (-0.2, 3.7) {BLE / MQTT};
\node[inputlbl] (lblFil) at (-0.2, 2.9) {Fichier XML};
\node[inputlbl] (lblRs)  at (-0.2, 2.1) {RS-232};

% =====================================================================
% Couche 1 : Adapter (teal pour distinguer la frontiere d'entree)
% =====================================================================
\node[altlayer] (adapter) at (1.4, 3.3) {Adapter\\Layer};

% Fleches d'entree
\draw[flowarrow] (lblTcp.east) -- (adapter.west |- lblTcp.east);
\draw[flowarrow] (lblBle.east) -- (adapter.west |- lblBle.east);
\draw[flowarrow] (lblFil.east) -- (adapter.west |- lblFil.east);
\draw[flowarrow] (lblRs.east)  -- (adapter.west |- lblRs.east);

% =====================================================================
% Couche 2 : Parser
% =====================================================================
\node[layer] (parser) at (4.6, 3.3) {Parser\\Layer};
\draw[flowarrow] (adapter.east) -- (parser.west)
  node[midway, above, font=\scriptsize\itshape, color=bridgegray, yshift=2pt]
    {trame brute};

% =====================================================================
% Couche 3 : Mapper
% =====================================================================
\node[layer] (mapper) at (7.8, 3.3) {Mapper\\Layer};
\draw[flowarrow] (parser.east) -- (mapper.west)
  node[midway, above, font=\scriptsize\itshape, color=bridgegray, yshift=2pt]
    {objet};

% =====================================================================
% Couche 4 : Transport
% =====================================================================
\node[altlayer] (transp) at (11.0, 3.3) {Transport\\Layer};
\draw[flowarrow] (mapper.east) -- (transp.west)
  node[midway, above, font=\scriptsize\itshape, color=bridgegray, yshift=2pt]
    {FHIR};

% =====================================================================
% Sortie HAPI
% =====================================================================
\node[outbox] (hapi) at (13.9, 3.3) {HAPI FHIR};
\draw[flowarrow] (transp.east) -- (hapi.west)
  node[midway, above, font=\scriptsize\itshape, color=bridgegray, yshift=2pt]
    {mTLS};

% =====================================================================
% Annotations en bas (bien sous chaque couche, espace suffisant)
% =====================================================================
\node[caption] at (1.4, 0.7) {drivers\\d'entree};
\node[caption] at (4.6, 0.7) {profils JSON\\decodage};
\node[caption] at (7.8, 0.7) {LOINC\\Observation FHIR};
\node[caption] at (11.0, 0.7) {retry, fallback\\file SQLite};

% =====================================================================
% Bandeau plugins (au-dessus, couvrant les 3 couches extensibles)
% =====================================================================
\node[draw=bridgegray, dashed, fill=bridgegray!8, rounded corners=2pt,
      minimum width=8.4cm, minimum height=0.7cm,
      anchor=south, font=\scriptsize\itshape\sffamily, color=bridgegray]
  at (4.6, 6.0)
  {Plugins communautaires : extensions Adapter, Parser, Mapper};

% Petites fleches verticales du bandeau vers les 3 couches
\foreach \x in {1.4, 4.6, 7.8} {
  \draw[->, dashed, color=bridgegray, thin] (\x, 5.95) -- (\x, 5.45);
}

\end{tikzpicture}
}
  \caption{Pipeline fonctionnel à 4 couches, avec annotations des formats à chaque interface et zone d'extension par plugins.}
  \label{fig:vue-fonctionnelle}
\end{figure}

\subsection{Vue applicative}

Six composants logiciels constituent la solution.

\begin{table}[h]
  \centering
  \caption{Composants applicatifs.}
  \label{tab:composants-app}
  \begin{tabularx}{\linewidth}{lX}
    \toprule
    \bridgeheader Composant & Responsabilité \\
    \midrule
    Adapter Layer & Drivers d'entrée par canal (TCP, MQTT, file-watcher, série). \\
    Parser Layer & Décodage des trames et fichiers selon un profil JSON par dispositif. \\
    Mapper Layer & Codification LOINC et construction de ressources FHIR R4. \\
    Transport Layer & Émission vers HAPI avec retry, fallback SQLite, mTLS optionnel. \\
    Dashboard & Configuration des profils et supervision OT (FastAPI + HTMX). \\
    Plugin Manager & Chargement, validation et hot-reload des plugins déposés dans \texttt{plugins/}. \\
    \bottomrule
  \end{tabularx}
\end{table}

\subsection{Vue technique}

L'environnement d'exécution est un poste local avec Docker Compose. Six services tournent en réseau interne : quatre simulateurs, le cœur Python, le serveur HAPI FHIR. Le dashboard est exposé sur \texttt{localhost:8080}. Aucun accès Internet n'est requis pendant la démo.

\begin{figure}[h]
  \centering
  \resizebox{\linewidth}{!}{% =====================================================================
% figures/fig-topologie-docker.tex
% Topologie Docker Compose : 4 simulateurs + bridge + HAPI + MQTT broker.
% Layout en 4 niveaux verticaux : sims, bus mqtt+volume, bridge+dashboard, hapi
% =====================================================================
\begin{tikzpicture}[
  font=\sffamily,
  every node/.style={font=\small\sffamily},
  >={Stealth[length=2.5mm]},
  thick,
  sim/.style={
    rectangle, rounded corners=3pt,
    draw=bridgealert, fill=bridgealert!10, text=bridgealert,
    minimum width=2.6cm, minimum height=1.0cm,
    align=center, font=\small\sffamily
  },
  bus/.style={
    rectangle, rounded corners=3pt,
    draw=bridgeteal, fill=bridgeteal!12, text=bridgeteal,
    minimum width=2.6cm, minimum height=0.9cm,
    align=center, font=\small\sffamily
  },
  vol/.style={
    rectangle, rounded corners=3pt,
    draw=bridgegray, dashed, fill=bridgegray!8,
    minimum width=2.6cm, minimum height=0.9cm,
    align=center, font=\small\sffamily, text=bridgegray
  },
  bridgecore/.style={
    rectangle, rounded corners=4pt,
    draw=bridgeblue, fill=bridgeblue!12, text=bridgeblue,
    minimum width=6.6cm, minimum height=1.6cm,
    align=center, font=\bfseries\sffamily
  },
  dashbox/.style={
    rectangle, rounded corners=3pt,
    draw=bridgeteal, fill=bridgeteal!15, text=bridgeteal,
    minimum width=2.6cm, minimum height=1.6cm,
    align=center, font=\small\bfseries\sffamily
  },
  hapi/.style={
    rectangle, rounded corners=3pt,
    draw=bridgegreen, fill=bridgegreen!15, text=bridgegreen,
    minimum width=3.0cm, minimum height=1.0cm,
    align=center, font=\bfseries\sffamily
  },
  protolbl/.style={
    font=\scriptsize\itshape\sffamily, color=bridgegray, inner sep=1pt
  },
  flowarrow/.style={-{Stealth[length=3mm]}, thick, draw=bridgeblue!80}
]

% =====================================================================
% Cadre du reseau Docker
% =====================================================================
\node[draw=bridgeblue, dashed, thick, fill=bridgesoft,
      rounded corners=4pt,
      minimum width=15.6cm, minimum height=10.0cm,
      anchor=south west,
      label={[anchor=north west, font=\footnotesize\bfseries,
              color=bridgeblue, xshift=8pt, yshift=-6pt]
              north west: Reseau Docker bridge-demo}]
      (net) at (0, 0) {};

% =====================================================================
% Niveau 1 : 4 simulateurs alignes en haut
% =====================================================================
\node[sim] (sim1) at (2.4, 8.6) {sim-monitor\\Philips};
\node[sim] (sim2) at (5.6, 8.6) {sim-glucometer\\Masimo};
\node[sim] (sim3) at (8.8, 8.6) {sim-spirometer\\Welch Allyn};
\node[sim] (sim4) at (12.0, 8.6) {sim-pump\\Drager};

% Etiquettes de protocole sous chaque sim, sans chevaucher
\node[protolbl] at (2.4, 7.85) {TCP HL7v2};
\node[protolbl] at (5.6, 7.85) {MQTT JSON};
\node[protolbl] at (8.8, 7.85) {fichier XML};
\node[protolbl] at (12.0, 7.85) {RS-232};

% =====================================================================
% Niveau 2 : Bus MQTT et volume partage
% =====================================================================
\node[bus] (mqtt) at (5.6, 6.3) {MQTT broker\\eclipse-mosquitto};
\node[vol] (vol)  at (8.8, 6.3) {Volume \texttt{drop/}};

% =====================================================================
% Niveau 3 : Bridge core (centre) et Dashboard (droite)
% =====================================================================
\node[bridgecore] (bridge) at (5.6, 4.0) {IoT Edge Bridge\\Python + FastAPI};
\node[font=\scriptsize\itshape, color=bridgegray, anchor=north]
  at (5.6, 3.4) {Adapter . Parser . Mapper . Transport};

\node[dashbox] (dashb) at (12.0, 4.0) {Dashboard\\:8080};

% =====================================================================
% Niveau 4 : HAPI FHIR Server
% =====================================================================
\node[hapi] (hapi) at (5.6, 1.4) {HAPI FHIR Server\\:8081};

% =====================================================================
% Acteur exterieur (auteur)
% =====================================================================
\node[anchor=west, font=\footnotesize\itshape, color=bridgegray]
  (user) at (16.0, 4.0) {Auteur};

% =====================================================================
% Fleches : simulateurs vers leurs cibles
% =====================================================================
% sim1 (TCP) directement vers Bridge
\draw[flowarrow] (sim1.south) -- ++(0, -0.3)
                 -| ([xshift=-2.0cm]bridge.north);

% sim2 (MQTT) vers MQTT broker
\draw[flowarrow] (sim2.south) -- (mqtt.north);

% sim3 (fichier) vers volume
\draw[flowarrow] (sim3.south) -- (vol.north);

% sim4 (RS-232) directement vers Bridge
\draw[flowarrow] (sim4.south) -- ++(0, -0.3)
                 -| ([xshift=2.0cm]bridge.north);

% MQTT et volume vers Bridge
\draw[flowarrow] (mqtt.south)  -- ([xshift=-0.7cm]bridge.north);
\draw[flowarrow] (vol.south)   -- ([xshift=0.7cm]bridge.north);

% Bridge vers HAPI
\draw[flowarrow] (bridge.south) -- (hapi.north)
  node[midway, right=2pt, font=\scriptsize\itshape, color=bridgegray] {POST FHIR};

% Bridge vers Dashboard
\draw[flowarrow] (bridge.east) -- (dashb.west)
  node[midway, above, font=\scriptsize\itshape, color=bridgegray] {WebSocket};

% =====================================================================
% Acces auteur vers Dashboard (en pointille, hors reseau Docker)
% =====================================================================
\draw[->, dashed, thick, color=bridgegray]
  (user.west) -- (dashb.east)
  node[midway, above, font=\scriptsize\itshape, color=bridgegray] {HTTP};

\end{tikzpicture}
}
  \caption{Vue technique : topologie Docker Compose de la démo, 4 simulateurs + bridge + HAPI FHIR + MQTT broker.}
  \label{fig:vue-technique}
\end{figure}
